\documentclass[10pt,a4paper]{amsart}
\usepackage{amsmath,amssymb,amsthm,graphicx}
\usepackage{verbatim}
\newenvironment{reference}{\comment}{\endcomment}
%\newenvironment{reference}{}{}
\newenvironment{slogan}{\comment}{\endcomment}
\newenvironment{history}{\comment}{\endcomment}

% For commutative diagrams we use Xy-pic
\usepackage[all]{xy}

% We use 2cell for 2-commutative diagrams.
\xyoption{2cell}
% Theorem environments.
%
\theoremstyle{plain}
\newtheorem{theorem}[subsection]{Theorem}
\newtheorem{proposition}[subsection]{Proposition}
\newtheorem{lemma}[subsection]{Lemma}

\theoremstyle{definition}
\newtheorem{definition}[subsection]{Definition}
\newtheorem{example}[subsection]{Example}
\newtheorem{exercise}[subsection]{Exercise}
\newtheorem{situation}[subsection]{Situation}

\theoremstyle{remark}
\newtheorem{remark}[subsection]{Remark}
\newtheorem{remarks}[subsection]{Remarks}

\numberwithin{equation}{subsection}

% Macros
%
\def\lim{\mathop{\mathrm{lim}}\nolimits}
\def\colim{\mathop{\mathrm{colim}}\nolimits}
\def\Spec{\mathop{\mathrm{Spec}}}
\def\Hom{\mathop{\mathrm{Hom}}\nolimits}
\def\Ext{\mathop{\mathrm{Ext}}\nolimits}
\def\SheafHom{\mathop{\mathcal{H}\!\mathit{om}}\nolimits}
\def\SheafExt{\mathop{\mathcal{E}\!\mathit{xt}}\nolimits}
\def\Sch{\mathit{Sch}}
\def\Mor{\mathop{\mathrm{Mor}}\nolimits}
\def\Ob{\mathop{\mathrm{Ob}}\nolimits}
\def\Sh{\mathop{\mathit{Sh}}\nolimits}
\def\NL{\mathop{N\!L}\nolimits}
\def\CH{\mathop{\mathrm{CH}}\nolimits}
\def\proetale{{pro\text{-}\acute{e}tale}}
\def\etale{{\acute{e}tale}}
\def\QCoh{\mathit{QCoh}}
\def\Ker{\mathop{\mathrm{Ker}}}
\def\Im{\mathop{\mathrm{Im}}}
\def\Coker{\mathop{\mathrm{Coker}}}
\def\Coim{\mathop{\mathrm{Coim}}}

% Boxtimes
%
\DeclareMathSymbol{\boxtimes}{\mathbin}{AMSa}{"02}

%
% Macros for moduli stacks/spaces
%
\def\QCohstack{\mathcal{QC}\!\mathit{oh}}
\def\Cohstack{\mathcal{C}\!\mathit{oh}}
\def\Spacesstack{\mathcal{S}\!\mathit{paces}}
\def\Quotfunctor{\mathrm{Quot}}
\def\Hilbfunctor{\mathrm{Hilb}}
\def\Curvesstack{\mathcal{C}\!\mathit{urves}}
\def\Polarizedstack{\mathcal{P}\!\mathit{olarized}}
\def\Complexesstack{\mathcal{C}\!\mathit{omplexes}}
% \Pic is the operator that assigns to X its picard group, usage \Pic(X)
% \Picardstack_{X/B} denotes the Picard stack of X over B
% \Picardfunctor_{X/B} denotes the Picard functor of X over B
\def\Pic{\mathop{\mathrm{Pic}}\nolimits}
\def\Picardstack{\mathcal{P}\!\mathit{ic}}
\def\Picardfunctor{\mathrm{Pic}}
\def\Deformationcategory{\mathcal{D}\!\mathit{ef}}
\usepackage{cleveref}
\crefname{originallemma}{original lemma}{original lemmas}
\Crefname{originallemma}{Original lemma}{Original lemmas}
\makeatletter
\newlabel{lemma-morphism-ringed-etale-topoi-affines}{{04I6}{source}{}{sourcepointer-1}{}}
\newlabel{lemma-morphism-ringed-etale-topoi-affines@cref}{{[originallemma][0][]04I6}{[1][0][]source}{}{}{}}
\makeatother
\crefname{originallemma}{original lemma}{original lemmas}
\Crefname{originallemma}{Original lemma}{Original lemmas}
\makeatletter
\newlabel{lemma-2-morphism}{{04IJ}{source}{}{sourcepointer-2}{}}
\newlabel{lemma-2-morphism@cref}{{[originallemma][0][]04IJ}{[1][0][]source}{}{}{}}
\makeatother
\crefname{originallemma}{original lemma}{original lemmas}
\Crefname{originallemma}{Original lemma}{Original lemmas}
\makeatletter
\newlabel{lemma-morphism-locally-ringed}{{04I5}{source}{}{sourcepointer-3}{}}
\newlabel{lemma-morphism-locally-ringed@cref}{{[originallemma][0][]04I5}{[1][0][]source}{}{}{}}
\makeatother
\crefname{originallemma}{original lemma}{original lemmas}
\Crefname{originallemma}{Original lemma}{Original lemmas}
\makeatletter
\newlabel{sites-modules-lemma-composition-morphisms-locally-ringed-topoi}{{04IG}{source}{}{sourcepointer-4}{}}
\newlabel{sites-modules-lemma-composition-morphisms-locally-ringed-topoi@cref}{{[originallemma][0][]04IG}{[1][0][]source}{}{}{}}
\makeatother
\crefname{originallemma}{original lemma}{original lemmas}
\Crefname{originallemma}{Original lemma}{Original lemmas}
\makeatletter
\newlabel{lemma-faithful}{{04LW}{source}{}{sourcepointer-5}{}}
\newlabel{lemma-faithful@cref}{{[originallemma][0][]04LW}{[1][0][]source}{}{}{}}
\makeatother
\crefname{originaldefinition}{original definition}{original definitions}
\Crefname{originaldefinition}{Original definition}{Original definitions}
\makeatletter
\newlabel{categories-definition-2-category}{{003H}{source}{}{sourcepointer-6}{}}
\newlabel{categories-definition-2-category@cref}{{[originaldefinition][0][]003H}{[1][0][]source}{}{}{}}
\makeatother
\crefname{originalsection}{original section}{original sections}
\Crefname{originalsection}{Original section}{Original sections}
\makeatletter
\newlabel{sites-section-2-category}{{04I9}{source}{}{sourcepointer-7}{}}
\newlabel{sites-section-2-category@cref}{{[originalsection][0][]04I9}{[1][0][]source}{}{}{}}
\makeatother
\crefname{originallemma}{original lemma}{original lemmas}
\Crefname{originallemma}{Original lemma}{Original lemmas}
\makeatletter
\newlabel{lemma-support-subsheaf-final}{{04HV}{source}{}{sourcepointer-8}{}}
\newlabel{lemma-support-subsheaf-final@cref}{{[originallemma][0][]04HV}{[1][0][]source}{}{}{}}
\makeatother
\crefname{originallemma}{original lemma}{original lemmas}
\Crefname{originallemma}{Original lemma}{Original lemmas}
\makeatletter
\newlabel{sites-modules-lemma-localize-morphism-locally-ringed-topoi}{{04IH}{source}{}{sourcepointer-9}{}}
\newlabel{sites-modules-lemma-localize-morphism-locally-ringed-topoi@cref}{{[originallemma][0][]04IH}{[1][0][]source}{}{}{}}
\makeatother
\title{Morphisms of locally ringed topoi on small etale topoi come from morphisms of schemes}
\author{The Stacks Project Authors}
\date{stacks-project:etale-cohomology.tex:local-snapshot-20261011; GNU FDL 1.2 or later}
\begin{document}
\maketitle
\noindent\textit{Editorial note.} All delivered mathematical text is the author's own, copied from the pinned Stacks Project chapter file etale-cohomology.tex; no statement or proof step was rewritten, reordered or removed.\par\medskip
\noindent\textit{Editorial note.} The proof is the author's complete construction: existence in the affine case, the global gluing, the compatibility of the 2-isomorphisms and uniqueness are all written out. The closing sentence 'Some details omitted' concerns routine compatibility checks in the gluing and is preserved verbatim in the counted span.\par\medskip
\noindent\textit{Editorial note.} Every labelled reference occurring inside the counted spans is carried in assembly.external\_references as a non-load-bearing navigation pointer (same-chapter ones at their definition site, cross-chapter ones at the invocation inside the counted span), so no load-bearing mathematics is replaced by a stub. No bibliography citation occurs inside the counted spans.\par\medskip
\noindent\textit{Editorial note.} The checked chapter list for same-name collisions is the set of stacks chapters staged in this repository: chow, cohomology, divisors, etale-cohomology, flat, more-algebra, more-morphisms, morphisms, spaces-more-morphisms, varieties together with this card's chapter etale-cohomology.\par\medskip
\noindent\textit{Editorial note.} The preamble is the chapter's own native preamble excerpt already used for this source: the Xy-pic setup (native preamble lines 14-18) and the native theorem declarations with the mathematical macros (native preamble lines 154-217).\par\medskip
\noindent\textit{Allowed standard prerequisites (editorial).}\begin{itemize}
\item The affine case of the theorem. Same chapter, Lemma lemma-morphism-ringed-etale-topoi-affines, whose proof is adjusted to the global setting. Supplies the local existence statement that is glued.
\item Rigidity of morphisms of locally ringed topoi. Same chapter, Lemma lemma-2-morphism. Makes the local morphisms agree on overlaps and the 2-isomorphisms compatible.
\item Localization morphisms of locally ringed topoi attached to open immersions. Same chapter, Lemma lemma-morphism-locally-ringed. Supplies the morphisms along which the local statements are composed.
\end{itemize}\medskip
\ifdefined\hypertarget\hypertarget{sourcepointer-1}{}\fi
\noindent\textit{Source reference.} Original lemma 04I6: sources/stacks-etale-cohomology/pinned-original-source.tex, source lines 5347--5347. Navigation pointer invoked at proof line 5576; the counted span contains the pointer only and the cited result is not reproduced inside it, so the pointer is registered as a non-load-bearing navigation stub in the established form for Stacks material. No bibliography citation is involved.\par\medskip
\ifdefined\hypertarget\hypertarget{sourcepointer-2}{}\fi
\noindent\textit{Source reference.} Original lemma 04IJ: sources/stacks-etale-cohomology/pinned-original-source.tex, source lines 5227--5227. Navigation pointer invoked at proof line 5580; the counted span contains the pointer only and the cited result is not reproduced inside it, so the pointer is registered as a non-load-bearing navigation stub in the established form for Stacks material. No bibliography citation is involved.\par\medskip
\ifdefined\hypertarget\hypertarget{sourcepointer-3}{}\fi
\noindent\textit{Source reference.} Original lemma 04I5: sources/stacks-etale-cohomology/pinned-original-source.tex, source lines 5178--5178. Navigation pointer invoked at proof line 5592; the counted span contains the pointer only and the cited result is not reproduced inside it, so the pointer is registered as a non-load-bearing navigation stub in the established form for Stacks material. No bibliography citation is involved.\par\medskip
\ifdefined\hypertarget\hypertarget{sourcepointer-4}{}\fi
\noindent\textit{Source reference.} Original lemma 04IG: sources/stacks-etale-cohomology/pinned-original-source.tex, source lines 5603--5603. Navigation pointer invoked at proof line 5603; the counted span contains the pointer only and the cited result is not reproduced inside it, so the pointer is registered as a non-load-bearing navigation stub in the established form for Stacks material. No bibliography citation is involved.\par\medskip
\ifdefined\hypertarget\hypertarget{sourcepointer-5}{}\fi
\noindent\textit{Source reference.} Original lemma 04LW: sources/stacks-etale-cohomology/pinned-original-source.tex, source lines 5298--5298. Navigation pointer invoked at proof line 5628; the counted span contains the pointer only and the cited result is not reproduced inside it, so the pointer is registered as a non-load-bearing navigation stub in the established form for Stacks material. No bibliography citation is involved.\par\medskip
\ifdefined\hypertarget\hypertarget{sourcepointer-6}{}\fi
\noindent\textit{Source reference.} Original definition 003H: sources/stacks-etale-cohomology/pinned-original-source.tex, source lines 5647--5647. Navigation pointer invoked at proof line 5647; the counted span contains the pointer only and the cited result is not reproduced inside it, so the pointer is registered as a non-load-bearing navigation stub in the established form for Stacks material. No bibliography citation is involved.\par\medskip
\ifdefined\hypertarget\hypertarget{sourcepointer-7}{}\fi
\noindent\textit{Source reference.} Original section 04I9: sources/stacks-etale-cohomology/pinned-original-source.tex, source lines 5649--5649. Navigation pointer invoked at proof line 5649; the counted span contains the pointer only and the cited result is not reproduced inside it, so the pointer is registered as a non-load-bearing navigation stub in the established form for Stacks material. No bibliography citation is involved.\par\medskip
\ifdefined\hypertarget\hypertarget{sourcepointer-8}{}\fi
\noindent\textit{Source reference.} Original lemma 04HV: sources/stacks-etale-cohomology/pinned-original-source.tex, source lines 3923--3923. Navigation pointer invoked at proof line 5667; the counted span contains the pointer only and the cited result is not reproduced inside it, so the pointer is registered as a non-load-bearing navigation stub in the established form for Stacks material. No bibliography citation is involved.\par\medskip
\ifdefined\hypertarget\hypertarget{sourcepointer-9}{}\fi
\noindent\textit{Source reference.} Original lemma 04IH: sources/stacks-etale-cohomology/pinned-original-source.tex, source lines 5670--5670. Navigation pointer invoked at proof line 5670; the counted span contains the pointer only and the cited result is not reproduced inside it, so the pointer is registered as a non-load-bearing navigation stub in the established form for Stacks material. No bibliography citation is involved.\par\medskip
\begin{theorem}
\label{theorem-fully-faithful}
Let $X$, $Y$ be schemes. Let
$$
(g, g^\#) :
(\Sh(X_\etale), \mathcal{O}_X)
\longrightarrow
(\Sh(Y_\etale), \mathcal{O}_Y)
$$
be a morphism of locally ringed topoi. Then there exists a
unique morphism of schemes $f : X \to Y$ such that
$(g, g^\#)$ is isomorphic to $(f_{small}, f_{small}^\sharp)$.
In other words, the construction
$$
\Sch \longrightarrow \textit{Locally ringed topoi},
\quad
X \longrightarrow (X_\etale, \mathcal{O}_X)
$$
is fully faithful (morphisms up to $2$-isomorphisms on the right hand side).
\end{theorem}

\begin{proof}
You can prove this theorem by carefully adjusting the arguments of
the proof of
Lemma \ref{lemma-morphism-ringed-etale-topoi-affines}
to the global setting. However, we want to indicate how we
can glue the result of that lemma to get a global morphism
due to the rigidity provided by the result of
Lemma \ref{lemma-2-morphism}.
Unfortunately, this is a bit messy.

\medskip\noindent
Let us prove existence when $Y$ is affine. In this case choose an
affine open covering $X = \bigcup U_i$. For each $i$ the inclusion
morphism $j_i : U_i \to X$ induces a morphism of locally ringed topoi
$(j_{i, small}, j_{i, small}^\sharp) :
(\Sh(U_{i, \etale}), \mathcal{O}_{U_i})
\to
(\Sh(X_\etale), \mathcal{O}_X)$
by
Lemma \ref{lemma-morphism-locally-ringed}.
We can compose this with $(g, g^\sharp)$ to obtain a morphism
of locally ringed topoi
$$
(g, g^\sharp) \circ (j_{i, small}, j_{i, small}^\sharp) :
(\Sh(U_{i, \etale}), \mathcal{O}_{U_i})
\to
(\Sh(Y_\etale), \mathcal{O}_Y)
$$
see
Modules on Sites,
Lemma \ref{sites-modules-lemma-composition-morphisms-locally-ringed-topoi}.
By
Lemma \ref{lemma-morphism-ringed-etale-topoi-affines}
there exists a unique morphism of schemes $f_i : U_i \to Y$
and a $2$-isomorphism
$$
t_i :
(f_{i, small}, f_{i, small}^\sharp)
\longrightarrow
(g, g^\sharp) \circ (j_{i, small}, j_{i, small}^\sharp).
$$
Set $U_{i, i'} = U_i \cap U_{i'}$, and denote $j_{i, i'} : U_{i, i'} \to U_i$
the inclusion morphism. Since we have
$j_i \circ j_{i, i'} = j_{i'} \circ j_{i', i}$
we see that
\begin{align*}
(g, g^\sharp) \circ
(j_{i, small}, j_{i, small}^\sharp) \circ
(j_{i, i', small}, j_{i, i', small}^\sharp)
= \\
(g, g^\sharp) \circ
(j_{i', small}, j_{i', small}^\sharp) \circ
(j_{i', i, small}, j_{i', i, small}^\sharp)
\end{align*}
Hence by uniqueness (see
Lemma \ref{lemma-faithful})
we conclude that
$f_i \circ j_{i, i'} = f_{i'} \circ j_{i', i}$, in other words the
morphisms of schemes $f_i = f \circ j_i$ are the restrictions of a
global morphism of schemes $f : X \to Y$. Consider the diagram
of $2$-isomorphisms (where we drop the components ${}^\sharp$ to ease the
notation)
$$
\xymatrix{
g \circ j_{i, small} \circ j_{i, i', small}
\ar[rr]^{t_i \star \text{id}_{j_{i, i', small}}}
\ar@{=}[d] & &
f_{small} \circ j_{i, small} \circ j_{i, i', small} \ar@{=}[d] \\
g \circ j_{i', small} \circ j_{i', i, small}
\ar[rr]^{t_{i'} \star \text{id}_{j_{i', i, small}}} & &
f_{small} \circ j_{i', small} \circ j_{i', i, small}
}
$$
The notation $\star$ indicates horizontal composition, see
Categories, Definition \ref{categories-definition-2-category}
in general and
Sites, Section \ref{sites-section-2-category}
for our particular case. By the result of
Lemma \ref{lemma-2-morphism}
this diagram commutes. Hence for any sheaf $\mathcal{G}$
on $Y_\etale$ the isomorphisms
$t_i : f_{small}^{-1}\mathcal{G}|_{U_i} \to g^{-1}\mathcal{G}|_{U_i}$
agree over $U_{i, i'}$ and we obtain a global isomorphism
$t : f_{small}^{-1}\mathcal{G} \to g^{-1}\mathcal{G}$.
It is clear that this isomorphism is functorial in $\mathcal{G}$
and is compatible with the maps $f_{small}^\sharp$ and $g^\sharp$
(because it is compatible with these maps locally).
This proves the theorem in case $Y$ is affine.

\medskip\noindent
In the general case, let $V \subset Y$ be an affine open.
Then $h_V$ is a subsheaf of the final sheaf $*$ on $Y_\etale$.
As $g$ is exact we see that $g^{-1}h_V$ is a subsheaf of the final
sheaf on $X_\etale$. Hence by
Lemma \ref{lemma-support-subsheaf-final}
there exists an open subscheme $W \subset X$ such that $g^{-1}h_V = h_W$. By
Modules on Sites,
Lemma \ref{sites-modules-lemma-localize-morphism-locally-ringed-topoi}
there exists a commutative diagram of morphisms of locally ringed
topoi
$$
\xymatrix{
(\Sh(W_\etale), \mathcal{O}_W) \ar[r] \ar[d]_{g'} &
(\Sh(X_\etale), \mathcal{O}_X) \ar[d]^g \\
(\Sh(V_\etale), \mathcal{O}_V) \ar[r] &
(\Sh(Y_\etale), \mathcal{O}_Y)
}
$$
where the horizontal arrows are the localization morphisms
(induced by the inclusion morphisms $V \to Y$ and $W \to X$)
and where $g'$ is induced from $g$. By the result of the preceding
paragraph we obtain a morphism of schemes $f' : W \to V$ and
a $2$-isomorphism
$t : (f'_{small}, (f'_{small})^\sharp) \to (g', (g')^\sharp)$.
Exactly as before these morphisms $f'$ (for varying affine opens $V \subset Y$)
agree on overlaps by uniqueness, so we get a morphism $f : X \to Y$.
Moreover, the $2$-isomorphisms $t$ are compatible on overlaps by
Lemma \ref{lemma-2-morphism}
again and we obtain a global $2$-isomorphism
$(f_{small}, (f_{small})^\sharp) \to (g, (g)^\sharp)$.
as desired. Some details omitted.
\end{proof}

\end{document}
